%%%%%%%%%%%%%%%%%%%%%%%%%%%%%%%%%%%%%%%%
% CAPÍTULO - CONCLUSÃO

\chapter{COMENTÁRIOS FINAIS}
\label{chp:conc}

Este trabalho procurou mostrar as principais características de duas missões espaciais, a NanosatC-Br1 e a OSIRIS-REx, a primeira uma missão do INPE e UFSM, e a segunda uma missão da Nasa.

O NanosatC-Br1 utiliza um satélite relativamente pequeno, em uma órbita baixa de cerca de 600 km de altitude para coletar informações do SAMA e EEJ, e a missão tem três objetivos nos campos da ciência, tecnologia, e academia.

O OSIRIS-Rex pesa cerca de 2 toneladas, tem dimensões aproximadas de um cubo de 2,5 m de lado, e é uma missão para coletar amostras do asteroide Bennu.

Este trabalho procurou mostrar os aspectos principais das missões, como descrições, funcionamentos, e equipamentos utilizados. Como sugestão de trabalhos futuros, os resultados obtidos até o momento das missões poderiam ser investigados e apresentados.